\section{Exact elimination of finite-dimensional messages}
\label{sec:peps_message_expansion}

The next identities justify the preliminary message expansion in
\cite[Theorem~5.2]{OpenAI2026PolynomialPEPS}. They act on the
original transmitted register and arbitrary spectator memory. In particular,
no bound on the spectator dimension occurs.

\begin{definition}[A local bra and a message summand]
    \label{def:peps_local_message_summand}
    \lean{TNLean.PEPS.PairEffect.Word.localDiscard}
    \lean{TNLean.PEPS.PairEffect.Word.transferTerm}
    \leanok
    \uses{def:peps_party_layout}
    For a vector $u\in H$ at a party $a$, the local bra is
    $\langle u|\otimes I$, acting on $H\otimes K$ and discarding the $H$
    register. For parties $a,b$ and an orthonormal basis $(e_i)_{i\in I}$
    of $H$, the $i$th message summand first applies the local bra
    $\langle e_i|$ at $a$, then prepares $e_i$ at $b$.
\end{definition}

\begin{theorem}[Local message summands]
    \label{thm:peps_local_message_summands}
    \lean{TNLean.PEPS.PairEffect.Word.isAllowed_localDiscard}
    \lean{TNLean.PEPS.PairEffect.Word.eval_localDiscard_tmul}
    \lean{TNLean.PEPS.PairEffect.Word.isAllowed_transferTerm}
    \lean{TNLean.PEPS.PairEffect.Word.sources_transferTerm}
    \lean{TNLean.PEPS.PairEffect.Word.eval_transferTerm_tmul}
    \leanok
    A local bra defined by a normalized vector has norm at most one. Every message summand is an allowed product of local
    contractions, containing no pair source. On $x\otimes y$ its action is
    \begin{align}
        x\otimes y\longmapsto \langle e_i,x\rangle e_i\otimes y,
        \label{eq:peps_message_summand_action}
    \end{align}
    where the first output register is assigned to the receiver.
\end{theorem}
\begin{proof}
    \leanok
    \uses{def:peps_local_message_summand}
    The bra has norm at most $\|e_i\|=1$, and preparation of $e_i$ is an
    isometry. Their composition has the action~(\ref{eq:peps_message_summand_action}). Both operations
    are local at their prescribed parties.
\end{proof}

\begin{theorem}[The exact message identity]
    \label{thm:peps_exact_message_identity}
    \lean{TNLean.PEPS.PairEffect.Word.sum_eval_transferTerm}
    \leanok
    Summing the message summands over an orthonormal basis gives exactly the
    identity transfer from sender to receiver, tensored with the identity
    on the spectator memory:
    \begin{align}
        \sum_{i\in I} |e_i\rangle_b\langle e_i|_a\otimes I_K
        =I_{H\otimes K}.
        \label{eq:peps_exact_message_identity}
    \end{align}
    This includes zero-dimensional transmitted spaces.
\end{theorem}
\begin{proof}
    \leanok
    \uses{thm:peps_local_message_summands}
    The orthonormal-basis expansion gives $\sum_i\langle e_i,x\rangle e_i=x$.
    The asserted equality follows first on pure tensors and then by linearity.
\end{proof}

\begin{definition}[Prepending and composing monomials]
    \label{def:peps_compose_actual_monomials}
    \lean{TNLean.PEPS.PairEffect.PartyChain.prependWord}
    \lean{TNLean.PEPS.PairEffect.PartyChain.comp}
    \leanok
    Prepend a word to the initial word of a monomial. To compose two monomials
    with matching intermediate registers, retain the first monomial's effects
    in their original order and append the second monomial at its final word.
\end{definition}

\begin{theorem}[Operators and pair effects of composed monomials]
    \label{thm:peps_compose_actual_monomials}
    \lean{TNLean.PEPS.PairEffect.PartyChain.eval_prependWord}
    \lean{TNLean.PEPS.PairEffect.PartyChain.isAllowed_prependWord}
    \lean{TNLean.PEPS.PairEffect.PartyChain.effectCount_prependWord}
    \lean{TNLean.PEPS.PairEffect.PartyChain.eval_comp}
    \lean{TNLean.PEPS.PairEffect.PartyChain.isAllowed_comp}
    \lean{TNLean.PEPS.PairEffect.PartyChain.effectCount_comp}
    \leanok
    \uses{def:peps_compose_actual_monomials}
    An allowed monomial may be preceded by an allowed word containing no pair
    effects, without changing its number of pair effects. More generally,
    two monomials with matching intermediate registers have a chronological
    composition. Its operator is the composition of their operators, its
    pair-effect count is the sum of their counts, and it is allowed whenever
    both original monomials are allowed.
\end{theorem}
\begin{proof}
    \leanok
    \uses{def:peps_compose_actual_monomials}
    Prepend the word to the initial word of the monomial. For composition,
    retain the first monomial's successive effects in their original order
    and append the second monomial at its final word. Induction on the first
    monomial proves the operator identity, allowedness and the exact count.
\end{proof}

\begin{theorem}[The literal finite expansion of a message]
    \label{thm:peps_message_expansion_count}
    \lean{TNLean.PEPS.PairEffect.gate_toGate_eq_sum}
    \lean{TNLean.PEPS.PairEffect.PartyGate.transfer}
    \lean{TNLean.PEPS.PairEffect.PartyGate.length_transfer}
    \lean{TNLean.PEPS.PairEffect.PartyGate.isAllowed_effectCount_transfer}
    \lean{TNLean.PEPS.PairEffect.PartyGate.sum_norm_coeff_transfer}
    \lean{TNLean.PEPS.PairEffect.PartyGate.eval_transfer}
    \leanok
    The finite expansion consisting of the message summands, each with
    coefficient one, has exactly $\dim H$ terms and absolute coefficient sum
    $\dim H$. Every term is allowed and has no pair effect. The operator of
    this actual expansion is exactly the transmitted identity operator.
    For $\dim H=0$ the expansion is empty and represents the zero operator
    on the zero space.
\end{theorem}
\begin{proof}
    \leanok
    \uses{thm:peps_exact_message_identity}
    The expansion is indexed once by each basis vector. Its operator is the
    sum of the weighted monomial operators, so~(\ref{eq:peps_exact_message_identity})
    applies directly. Counting the unit coefficients proves both assertions
    about size.
\end{proof}

\begin{definition}[Composition of gate expansions]
    \label{def:peps_composed_gate_expansion}
    \lean{TNLean.PEPS.PairEffect.PartyGate.comp}
    \lean{TNLean.PEPS.PairEffect.PartyGate.mem_comp}
    \leanok
    \uses{def:peps_compose_actual_monomials}
    Let $L=((c_j,M_j))_j$ and $K=((d_k,N_k))_k$ be ordered gate
    expansions with matching intermediate memories. Their chronological
    composition lists $(c_jd_k,N_kM_j)$ in lexicographic order on $(j,k)$.
    Each composed occurrence comes from one occurrence in each original list.
\end{definition}

\begin{theorem}[Operator and resources of composed expansions]
    \label{thm:peps_composed_gate_resources}
    \lean{TNLean.PEPS.PairEffect.PartyGate.gate_comp}
    \lean{TNLean.PEPS.PairEffect.PartyGate.length_comp}
    \lean{TNLean.PEPS.PairEffect.PartyGate.sum_norm_comp}
    \lean{TNLean.PEPS.PairEffect.PartyGate.isAllowed_comp}
    \lean{TNLean.PEPS.PairEffect.PartyGate.effectCount_comp_le}
    \leanok
    \uses{def:peps_composed_gate_expansion}
    The composed list evaluates to the second gate operator after the
    first. Its length is $|L||K|$, and its absolute coefficient sum is
    $(\sum_j|c_j|)(\sum_k|d_k|)$. If the original monomials are allowed,
    so are their compositions. Bounds $r,s$ on their pair-effect counts
    give the bound $r+s$ on every composed occurrence.
\end{theorem}
\begin{proof}
    \leanok
    \uses{thm:peps_compose_actual_monomials}
    Distribute the two finite operator sums and use the exact operator
    identity for monomial composition. Counting the ordered pairs gives
    the length formula. Since $|c_jd_k|=|c_j||d_k|$, the coefficient
    sums multiply. The local contractions and pair effects remain in
    their original order, so allowedness is preserved and effect counts add.
\end{proof}
\begin{definition}[Monomials with finite-dimensional messages]
    \label{def:peps_message_monomial}
    \lean{TNLean.PEPS.PairEffect.MessageMonomial}
    \lean{TNLean.PEPS.PairEffect.MessageMonomial.eval}
    \lean{TNLean.PEPS.PairEffect.MessageMonomial.IsAllowed}
    \lean{TNLean.PEPS.PairEffect.MessageMonomial.messageDimensions}
    \lean{TNLean.PEPS.PairEffect.MessageMonomial.effectCount}
    \lean{TNLean.PEPS.PairEffect.MessageMonomial.expand}
    \leanok
    \uses{def:peps_party_monomial}
    A monomial with messages is a chronological composition of the
    existing source-and-effect monomials and finite-dimensional register
    transfers. A transfer is the identity on a register's Hilbert space,
    changing its owner from the sender to the receiver and leaving all
    spectators unchanged. The monomial is allowed when all its original
    source-and-effect segments are allowed.
    Record its message dimensions $d_1,\ldots,d_s$ in chronological order
    and let $r_M$ count its original pair effects. Choose an orthonormal
    basis of each message space and expand its identity as the sum of
    sender bras followed by receiver kets. Compose these finite lists in
    the original order to obtain an expansion $\mathcal E(M)$.
\end{definition}

\begin{theorem}[Exact expansion and resources of bounded messages]
    \label{thm:peps_message_expansion}
    \uses{def:peps_message_monomial}
    \lean{TNLean.PEPS.PairEffect.MessageMonomial.eval_expand}
    \lean{TNLean.PEPS.PairEffect.MessageMonomial.isAllowed_effectCount_expand}
    \lean{TNLean.PEPS.PairEffect.MessageMonomial.expansion_size}
    \leanok
    The expansion $\mathcal E(M)$ evaluates exactly to the original
    monomial. If $M$ is allowed, every term is an allowed existing
    monomial with exactly $r_M$ pair effects, and
    \begin{align}
        |\mathcal E(M)|
        =\sum_{(c,N)\in\mathcal E(M)}|c|
        =\prod_{j=1}^{s}d_j.
        \label{eq:peps_message_expansion_size}
    \end{align}
    These statements include zero-dimensional transfers and the empty
    list of message occurrences.
\end{theorem}
\begin{proof}
    \uses{thm:peps_message_expansion_count,thm:peps_composed_gate_resources}
    \leanok
    The basis expansion of a transfer has one unit-coefficient local
    monomial for each basis vector and introduces no pair effect.
    Composition multiplies the list lengths and absolute coefficient
    sums, adds the original effect counts, and preserves the evaluated
    operator. Induction over the actual monomial gives the assertions.
\end{proof}

\begin{theorem}[Polynomial cost of boundedly many messages]
    \label{thm:peps_message_expansion_polynomial}
    \lean{TNLean.PEPS.PairEffect.MessageMonomial.expansion_size_le}
    \lean{TNLean.PEPS.PairEffect.MessageMonomial.expansion_size_le_of_power_bounds}
    \leanok
    Suppose $M$ contains at most $m$ messages, each of dimension at most
    $D\geq1$. Both the number of expanded monomials and their absolute
    coefficient sum are at most $D^m$. In particular, if
    $d_j\leq C_D L^v$, where $C_D,L\geq1$ and $v$ is a nonnegative
    integer, both quantities are at most $C_D^mL^{vm}$.
    Thus boundedly many polynomial-dimensional messages change a gate
    expansion by only a polynomial factor, as required in Theorem~5.2,
    assumption~3, and the preliminary reduction in its proof.
\end{theorem}
\begin{proof}
    \leanok
    \uses{thm:peps_message_expansion}
    Bound each factor in~(\ref{eq:peps_message_expansion_size}) by $D$.
    Since $D\geq1$, replacing the actual message count by $m$ can only
    increase this upper bound. Substitute
    $D=C_D L^v$ for the final assertion.
\end{proof}
\section{Gates containing bounded messages}
\label{sec:peps_message_gate_expansion}

Let $H_a,H_b$ be the input and output memories of a fixed finite set of
participating parties. We use the monomials that permit local contractions,
normalized pair sources and effects, and finite-dimensional transfers, as in
% Source: 04-compression.tex, lines 32--43.
\cite[Proof of Theorem~5.2]{OpenAI2026PolynomialPEPS}.

\begin{definition}[Weighted message monomials]
    \label{def:peps_message_gate}
    \lean{TNLean.PEPS.PairEffect.MessageGate,
        TNLean.PEPS.PairEffect.MessageGate.eval}
    \leanok
    \uses{def:peps_message_monomial}
    A weighted gate list is a finite ordered list
    $L=((c_j,M_j))_{j=0}^{K-1}$, where each $M_j:H_a\to H_b$ is such a
    monomial and $c_j\in\C$. Its operator is $G_L=\sum_j c_jM_j$.
    Occurrences are retained separately when two monomials agree.
\end{definition}

\begin{definition}[Expanding all messages in a gate]
    \label{def:peps_message_gate_expand}
    \lean{TNLean.PEPS.PairEffect.MessageGate.expand,
        TNLean.PEPS.PairEffect.MessageGate.mem_expand}
    \leanok
    \uses{def:peps_message_gate}
    Write the coordinate expansion of the messages in $M_j$ as the ordered
    list $((a_{j,t},N_{j,t}))_t$. Replace the occurrence $(c_j,M_j)$ by
    $((c_ja_{j,t},N_{j,t}))_t$ and concatenate these lists in their original
    order. Thus every resulting occurrence has coefficient $c_ja_{j,t}$
    and monomial $N_{j,t}$ for an original occurrence $j$ and one of its
    coordinate choices $t$.
\end{definition}

\begin{theorem}[Operator and effect count under message expansion]
    \label{thm:peps_message_gate_eval}
    \lean{TNLean.PEPS.PairEffect.MessageGate.eval_expand,
        TNLean.PEPS.PairEffect.MessageGate.isAllowed_expand,
        TNLean.PEPS.PairEffect.MessageGate.effectCount_expand_le}
    \leanok
    \uses{def:peps_message_gate_expand}
    The expanded list has operator $G_L$. If every $M_j$ is allowed, so is
    every resulting monomial. If every $M_j$ is allowed and has at most $r$ pair effects,
    the same bound holds for every resulting monomial.
\end{theorem}
\begin{proof}
    \uses{thm:peps_message_expansion}
    \leanok
    Each coordinate expansion satisfies $\sum_t a_{j,t}N_{j,t}=M_j$.
    Hence $\sum_{j,t}c_ja_{j,t}N_{j,t}=\sum_jc_jM_j$. Its factors are the
    original allowed factors and the local coordinate maps of the transfers.
    The latter introduce no pair effects.
\end{proof}

\begin{theorem}[Term count and coefficient sum]
    \label{thm:peps_message_gate_size}
    \lean{TNLean.PEPS.PairEffect.MessageGate.expansion_size,
        TNLean.PEPS.PairEffect.MessageGate.expansion_size_le,
        TNLean.PEPS.PairEffect.MessageGate.expansion_size_le_of_coeff_le}
    \leanok
    \uses{def:peps_message_gate_expand}
    Let $d_{j,1},\ldots,d_{j,m_j}$ be the dimensions of the messages in
    $M_j$, and put $q_j=\prod_{t=1}^{m_j}d_{j,t}$. The number of resulting
    occurrences and their absolute coefficient sum are respectively
    $\sum_jq_j$ and $\sum_j|c_j|q_j$.
    If $|L|\le K$, $m_j\le m$, and $d_{j,t}\le D$ with $D\ge1$, then
    \begin{align}
        |\widetilde L| & \le KD^m,                    \\
        \sum_{(c,N)\in\widetilde L}|c|
                       & \le D^m\sum_{(c,M)\in L}|c|.
    \end{align}
    In particular, $|c_j|\le C$ with $C\ge0$ gives the bound $KCD^m$.
\end{theorem}
\begin{proof}
    \uses{thm:peps_message_expansion}
    \leanok
    A monomial with message dimensions $d_{j,t}$ has exactly $q_j$
    coordinate choices, and the absolute coefficient sum of its expansion
    is $q_j$. Multiplication by $c_j$ scales this sum by $|c_j|$.
    Summing proves the exact formulas. The inequalities follow from
    $q_j\le D^{m_j}\le D^m$ and $\sum_j|c_j|\le KC$.
\end{proof}

\begin{definition}[The corresponding original circuit gate]
    \label{def:peps_message_gate_circuit}
    \lean{TNLean.PEPS.PairEffect.MessageGate.toOriginalCircuit}
    \leanok
    \uses{thm:peps_message_gate_eval}
    Suppose the local participant set $C$ has $|C|\ne1$, its embedding in
    the global parties is $\iota:C\hookrightarrow P$, all monomials are
    allowed, and $\|G_L\|\le1$. Insert the expanded list as one nonprivate
    gate with the same participant embedding and the same spectator memory.
    Empty participant sets are included.
\end{definition}

\begin{theorem}[Operator and occurrences of the constructed gate]
    \label{thm:peps_message_gate_circuit_properties}
    \lean{TNLean.PEPS.PairEffect.MessageGate.eval_toOriginalCircuit,
        TNLean.PEPS.PairEffect.MessageGate.nonprivateCount_toOriginalCircuit,
        TNLean.PEPS.PairEffect.MessageGate.participants_toOriginalCircuit}
    \leanok
    \uses{def:peps_message_gate_circuit}
    The constructed gate's operator is the canonical relabelling of
    $G_L\otimes I$. It has one nonprivate occurrence, with participant
    set $\iota(C)$.
\end{theorem}
\begin{proof}
    \leanok
    \uses{thm:peps_message_gate_eval}
    The expanded list has the original operator, so adjoining the spectator
    identity and relabelling the participant registers gives the stated
    operator. Its norm bound follows from equality of the expanded and
    original operators. The construction inserts exactly one nonprivate
    occurrence with the prescribed participant embedding.
\end{proof}

\begin{theorem}[Resource bounds for the constructed gate]
    \label{thm:peps_message_gate_circuit_bounds}
    \lean{TNLean.PEPS.PairEffect.MessageGate.isExpansionBounded_toOriginalCircuit,
        TNLean.PEPS.PairEffect.MessageGate.isMonomialBounded_toOriginalCircuit}
    \leanok
    \uses{def:peps_message_gate_circuit}
    Under the dimension and message-count bounds of
    Theorem~\ref{thm:peps_message_gate_size}, if the original coefficient
    sum is at most $S$ and every original monomial has at most $r$ pair
    effects, the constructed gate has the same effect bound and coefficient
    sum at most $SD^m$. If the original list has at most $K$ terms, it has
    at most $\lceil KD^m\rceil$ terms.
\end{theorem}
\begin{proof}
    \leanok
    \uses{thm:peps_message_gate_eval,thm:peps_message_gate_size}
    Apply the effect-count preservation and the two size estimates to the
    actual expanded list used in the gate. The ceiling converts its real
    upper bound into an integer bound.
\end{proof}
